\documentclass[../main.tex]{subfiles}
\begin{document}

\chapter{The Translators and Their Works}\label{chap:translators}

\section{Transliteration}\label{sec:translit}

% A comment that must not appear in the output.
The school of \hunayn ibn \ishaq produced the \jaw, as reported by
\usaybia and \qifti in his \qiftitarikh. \ishaq's\footnote{Note the
possessive after a macro with \texttt{xspace}.} contemporary \thabit
ibn Qurrah, and \hubaysh al-A`sam, worked under al-Ma'm\={u}n.
\ali ibn Ri\d{d}w\={a}n wrote against \dimashqi.
Raw forms: al-\v{s}ay\U{h}, \d{h}ad\={\i}\b{t}, \={I}s\={a}, \.{g}ayr,
\kha\={a}\thaa{} and \Jim al\={\i}n\=us; \Fi al-\d{t}ibb.

\section{Languages}\label{sec:languages}

Galen's treatise \textgreek{\Peri αἱρέσεων τοῖς εἰσαγομένοις} (\textit{De sectis})
was known in Arabic as \textarabic{كتاب الفرق} (\textit{Kit\={a}b al-Firaq}),
see \cref{sec:translit} and \cref{tab:witnesses} on p.~\pageref{tab:witnesses}.
The Hebrew version reads \texthebrew{ספר הכתות} in one witness.\footnote{A
footnote with Arabic: \textarabic{قال جالينوس} (\textit{q\={a}la
\v{G}\={a}l\={\i}n\=us}), and Greek \foreignlanguage{greek}{Γαληνός}.}

\begin{Arabic}
قال حنين بن إسحاق: إن كتب جالينوس في الطب كثيرة (١٦ كتابا)، وقد ترجمتها إلى العربية.
\end{Arabic}

An English paragraph follows the Arabic block, with an inline RTL phrase
at the end of a sentence: the term is \textarabic{جوامع}.

\subsection{Citations}

As \textcite{Ullmann1970} showed, the summaries circulated
widely.\footnote{\cite[59]{Ullmann1970}; see also
\cite[cf.][123--125]{SavageSmith2002}.} Later scholarship
\parencite[e.g.][12]{Koetschet2019} confirms this
\autocite{Pormann2026}.\footcite[\textarabic{ص}~٣٤]{Ibn-Abi-Usaybia}
The Paris manuscript \cite{MS-Paris-2860} preserves the text.

\subsubsection{Draft markup}

Some text is \hl{highlighted} and some is \textcolor{red}{red}.
An undefined macro \unknownmacro{with argument} must survive.

\begin{itemize}
\item First item with \emph{emphasis}.
\item Second item with \textbf{bold}.
\end{itemize}

\section{Tables, Mathematics and Figures}

\begin{table}[ht]
\centering
\caption{Witnesses of the \jaw}\label{tab:witnesses}
\begin{tabular}{lll}
\toprule
Siglum & Library & Shelfmark \\
\midrule
\multirow{2}{*}{P} & \multicolumn{2}{c}{Paris, BnF} \\
 & arabe & 2860 \\
I & \textarabic{إسطنبول} & Aya Sofya 3593 \\
\bottomrule
\end{tabular}
\end{table}

\begin{longtable}{p{3cm}p{8cm}}
\caption{Long list of terms}\label{tab:terms}\\
\toprule
Greek & Arabic \\
\midrule
\endhead
\textgreek{αἵρεσις} & \textarabic{فرقة} \\
\textgreek{ἐμπειρία} & \textarabic{تجربة} \\
\bottomrule
\end{longtable}

Inline mathematics $a^2 + b^2 = c^2$ and a display:
\begin{equation}\label{eq:humours}
H = \sum_{i=1}^{4} \frac{q_i}{\sqrt{t}}
\end{equation}
as in \cref{eq:humours}.

\begin{figure}[ht]
\centering
\includegraphics[width=0.5\textwidth]{figures/stemma.pdf}
\caption{Stemma codicum}\label{fig:stemma}
\end{figure}

\begin{figure}[ht]
\centering
\begin{tikzcd}
\omega \arrow[r] \arrow[d] & P \\
\alpha \arrow[r] & I
\end{tikzcd}
\caption{Transmission diagram (tikz-cd)}\label{fig:tikzcd}
\end{figure}

\end{document}
